\section{幻灯片精华}

\subsection{核心概念图}

\begin{figure}[H]
\centering
\includegraphics[width=0.92\textwidth,page=3]{05_cs224r_offpolicy_actor_critic_2025.pdf}
\caption{Lecture slide 重述 off-policy actor-critic 数据流与 surrogate objective。}
\end{figure}
\clearpage

\subsection{PPO vs SAC 结构对比}

\begin{figure}[H]
\centering
\includegraphics[width=0.92\textwidth,page=5]{05_cs224r_offpolicy_actor_critic_2025.pdf}
\caption{Slide 中对比图展示 PPO 的 clipping objective 与 SAC 的 entropy regularization。}
\end{figure}
\clearpage

\subsection{Q 函数稳定性案例}

\begin{figure}[H]
\centering
\includegraphics[width=0.92\textwidth,page=6]{05_cs224r_offpolicy_actor_critic_2025.pdf}
\caption{Slide 展示 Q 函数爆炸 vs smoothing 的实证曲线。}
\end{figure}
\clearpage

\subsection{Replay Buffer 与 Sim2Real}

\begin{figure}[H]
\centering
\includegraphics[width=0.92\textwidth,page=8]{05_cs224r_offpolicy_actor_critic_2025.pdf}
\caption{Slide 演示 prioritized replay 与 sim-to-real pipeline。}
\end{figure}
\clearpage

\subsection{策略蒸馏流程图}

\begin{figure}[H]
\centering
\includegraphics[width=0.92\textwidth,page=9]{05_cs224r_offpolicy_actor_critic_2025.pdf}
\caption{Slide 展示 Sim2Real 与 policy distillation 的交互流程。}
\end{figure}
\clearpage

\subsection{性能雷达}

\begin{figure}[H]
\centering
\includegraphics[width=0.92\textwidth,page=10]{05_cs224r_offpolicy_actor_critic_2025.pdf}
\caption{Slide 展示 PPO/SAC 在不同维度（data efficiency, stability, deployment cost）上的 radar chart。}
\end{figure}
\clearpage

\subsection{模拟-真实迁移示意}

\begin{figure}[H]
\centering
\includegraphics[width=0.92\textwidth,page=12]{05_cs224r_offpolicy_actor_critic_2025.pdf}
\caption{Slide 概括 sim-to-real pipeline 中的 domain randomization 与 fallback 机制。}
\end{figure}
\clearpage

\subsection{训练与部署工作流}

\begin{figure}[H]
\centering
\includegraphics[width=0.92\textwidth,page=13]{05_cs224r_offpolicy_actor_critic_2025.pdf}
\caption{Slide 给出了训练实验、日志、评估、部署的闭环流程图。}
\end{figure}
\clearpage

\subsection{调试与日志示例}

\begin{figure}[H]
\centering
\includegraphics[width=0.92\textwidth,page=14]{05_cs224r_offpolicy_actor_critic_2025.pdf}
\caption{Slide 展示常用的 debugging metrics（KL divergence、return variance、entropy）。}
\end{figure}
\clearpage

\subsection{Overfit 与 Regularization 诊断}

\begin{figure}[H]
\centering
\includegraphics[width=0.92\textwidth,page=15]{05_cs224r_offpolicy_actor_critic_2025.pdf}
\caption{Slide 演示 Overfit 情况的 reward 曲线与 regularization 对比。}
\end{figure}
\clearpage

\subsection{使用 entropy 进行监控}

\begin{figure}[H]
\centering
\includegraphics[width=0.92\textwidth,page=16]{05_cs224r_offpolicy_actor_critic_2025.pdf}
\caption{Slide 展示 entropy decay 曲线与 policy action distribution。}
\end{figure}
\clearpage

\subsection{部署后的长期监控}

\begin{figure}[H]
\centering
\includegraphics[width=0.92\textwidth,page=17]{05_cs224r_offpolicy_actor_critic_2025.pdf}
\caption{Slide 提示在部署后需要监控的 drift 以及安全阈值。}
\end{figure}
\clearpage

\subsection{本章小结}

幻灯片中的可视化帮助总结 actor-critic 数据流、PPO/SAC 的权衡、Q function 的平滑处理与回放架构在 Sim2Real 中的应用，为后续工程提供直观参考。

